% SHARD-ID: CA-40-QUBIT-1D-RESIDUAL-COMPUTER-ALGEBRA
% SHARD-TITLE: Computer Algebra for One-Dimensional Qubit Residuals
% SHARD-SUMMARY: Records the Julia helpers that reconstruct Pauli operators, currents, coboundaries, and boost residual densities.
% SHARD-SUMMARY: The tests check nontrivial invariants: current collapse for symmetric onsite terms, nonzero interacting currents, and coboundary operator reconstruction.
% SHARD-SUMMARY: The code is a finite residual engine, not yet an SDP solver.
% SHARD-KEYWORDS: Julia, Pauli coefficients, residuals, coboundary, tests, qubit chain

\section{Computer Algebra for One-Dimensional Qubit Residuals}
\label{sec:qubit-1d-residual-computer-algebra}

\subsection{Status, Sources, and Dependencies}

\textbf{Status: checked.}  This shard documents the finite algebra engine
behind CA-35 and CA-39.

Dependencies: CA-35, CA-38, CA-39, CONVENTIONS.md (m), (n).

% Checked: src/QubitPauliLattice.jl, src/QubitPauliResiduals.jl, and
%   test/runtests.jl, testset "qubit Pauli nearest-neighbour current
%   obstructions".

\subsection{Operator and Coefficient Roundtrip}

The file \texttt{src/QubitPauliLattice.jl} implements
\[
  \sum_{\alpha_0,\ldots,\alpha_{n-1}=0}^3
  c_{\alpha_0\ldots\alpha_{n-1}}
  \sigma^{\alpha_0}\otimes\cdots\otimes\sigma^{\alpha_{n-1}}
\]
for \(n=1,2,3\) and for general \(n\).  The inverse map uses
\[
  c_\alpha=2^{-n}\Tr(P_\alpha A).
\]
The tests assert coefficient roundtrips for two-site densities and for the
five-site residual representatives.  The code rejects arrays whose axes are
not all length \(4\), and rejects non-self-adjoint matrix inputs.

\subsection{Checked Residual Helpers}

The current helper returns
\[
  p=i[h_{12},h_{23}]
\]
and the direct coefficient helper returns exactly (39.1).  The conservation
helper
\[
  A=i[h_{12}+h_{23},p]
\]
is the minimal three-site density for \(i[H,P]\) after the Jacobi cancellation
of outer-overlap terms.

The file \texttt{src/QubitPauliResiduals.jl} adds
\[
  (D u)_{a_0\ldots a_n}
  =
  u_{a_0\ldots a_{n-1}}\delta_{a_n0}
  -
  \delta_{a_00}u_{a_1\ldots a_n}
\]
and
\[
  B=i[p_{123},h_{23}]+2i[p_{123},h_{34}].
\]
The test suite checks that the coboundary coefficients reconstruct the operator
identity
\[
  D u = u_{12}-u_{23}
\]
on a three-site window, including the fact that the scalar \(u_{00}I\) is in
the kernel.

\subsection{Sentinels}

The checked examples are deliberately small:

\begin{itemize}
\item A generic real one-site density embedded by the symmetric split has
  \(p=0\), \(A=0\), and \(B=0\), while the embedded Hamiltonian density itself
  is nonzero.  This is the concrete fully-local exclusion witness for CA-39.
\item A pure classical \(ZZ\) density also has \(p=0\).
\item A transverse-Ising-style density has nonzero \(p\) and nonzero boost
  residual density \(B\).  This only means it survives the first obstruction.
\item An asymmetric fake on-site split creates a nonzero current, but the
  conservation sentinel is larger than the current and catches the density-split
  artifact.
\end{itemize}

\subsection{Current Gap}

The code does not yet solve the polynomial system (39.1), (39.4), (39.5), and
does not build moment matrices.  Its role is narrower: it makes every displayed
coefficient tensor executable and gives mutation-friendly invariants before an
SDP exporter is added.

